% main.tex — Phase-11 paper master file
%
% Compile with:
%     latexmk -pdf paper/main.tex
% The build runs from the repository root; all \input and
% \includegraphics paths are relative to paper/.

\documentclass[10pt,letterpaper,twocolumn]{article}

\usepackage{preamble}

% BibLaTeX with alphabetic style per prompts/11_paper.md §Mandatory
% Style Rules #1.
\usepackage[backend=biber,style=alphabetic,sorting=ynt]{biblatex}
\addbibresource{references.bib}

\title{\project}

\author{%
  MCP Isolation Research Authors\\
  \texttt{\{authors\}@mcp-iso.research}%
}

\date{July 2026}

\begin{document}

\maketitle

\begin{abstract}
\noindent
The \mcp{} distributes isolation across eight implicit trust
boundaries without mandating enforcement at any of them. As
multi-tenant deployments emerge, cross-tenant leakage becomes a
first-class concern: we present the first systematic study of
\emph{cross-tenant} MCP isolation, complementing Chen et al.'s
single-tenant runtime study. We catalogue all eight boundaries
with STRIDE attack-surface rows (63 ticket IDs, four misuse
cases), ship an open attack library of 50 concrete classes
(48 per-boundary plus prompt injection and grammar fuzzing), and
introduce a reproducible framework that measures cross-tenant
leakage under controlled concurrency. Across four research
questions and 30 iterations per cell, cache attacks dominate the
vulnerable server's leakage volume (85.7\%, $z = 40.73$,
$p < 10^{-4}$); the composition of four protocol-layer defenses
(per-tenant tool registry, tenant-prefixed cache keys,
resource-path canonicalisation, audience-bound JWTs) reduces
leak rate to zero (sign test $p = 0.031$, Wilcoxon
$p = 0.031$); no single defense closes the gap alone. We close
with a layered defense blueprint and three concrete
spec-change recommendations for MCP-2. \emph{All experiments use
synthetic tenant data; no real user data, credentials, or LLM
outputs are involved. The vulnerable reference server runs
inside a Firecracker microVM.}
\end{abstract}

% Optional keywords line (USENIX uses \beginkeywords).
\begin{keywords}
Model Context Protocol, multi-tenant isolation, MCP security,
prompt injection, cross-tenant leakage, STRIDE, fuzzing,
defense-in-depth, capability-based security.
\end{keywords}

\section{Introduction}
% 01_introduction.tex

The \mcp{} (\mcp{}) was introduced in late 2024 as an open
standard for connecting LLM agents to external tools, resources,
and prompt templates~\cite{mcp_spec,anthropic_mcp_announce}. Its
adoption has been rapid: every major model provider now ships an
MCP client, and a public registry lists thousands of community
servers ranging from filesystem gateways to database frontends.
The protocol's appeal is its simplicity---a single JSON-RPC
envelope with three core methods (\texttt{tools/call},
\texttt{resources/read}, \texttt{prompts/get}) plus a thin
\texttt{initialize} handshake---but that simplicity comes at a
cost. MCP distributes \emph{isolation} across eight implicit
trust boundaries (transport, session, namespace, tool, resource,
memory, cache, authentication) without mandating concrete
enforcement at any of them. As deployments move from
single-tenant desktop setups to multi-tenant server fleets, the
question of \emph{cross-tenant isolation} becomes central: can
one tenant observe another's tool calls, results, cached state,
or session metadata? The MCP specification is silent on the
question, and mainstream reference implementations inherited the
silence.

\paragraph{Why this matters.}
A multi-tenant MCP server that serves multiple users, agents, or
organisations inherits the threat model of any multi-tenant
service: each tenant's data and computation must be confined to
that tenant. The recent empirical study by Chen et al.~(2026)
finds that runtime MCP servers routinely share stdio workers,
follow symlinks across tenant roots, and key caches without
tenant identity~\cite{chen_2026_mcp_security}. Concurrent work
argues that isolation is a first-class principle for LLM-agent
safety but provides neither a concrete boundary catalogue nor
empirical measurement~\cite{jing_2026_isolation}. The gap is
empirical and architectural: no prior work has systematically
catalogued MCP isolation failures across all eight boundaries,
mapped them to automated attack patterns, and measured their
leakage under a reproducible multi-tenant harness.

\paragraph{Our approach.}
We close the gap with a five-part study:

\begin{enumerate}
  \item \textbf{An isolation-boundary catalogue} for MCP. We
        enumerate all eight boundaries and tag each with STRIDE
        rows~\cite{stride_orig} and 63 forward ticket IDs
        ($A$-$\{TRN,SES,NSP,TOL,RES,MEM,CCH,AUT\}$-$\{nnn\}$).
  \item \textbf{A 14-concept security taxonomy} with explicit
        MCP-boundary bindings, covering prompt injection,
        confused deputy, capability-based access control,
        sandboxing, side-channel leakage, cache poisoning,
        memory poisoning, namespace collision, tool squatting,
        resource traversal, request smuggling, and authentication
        bypass.
  \item \textbf{A reproducible empirical framework} that measures
        cross-tenant leakage in MCP deployments under controlled
        multi-tenant concurrency, reporting
        \emph{leakage\_rate}, \emph{time\_to\_leak},
        \emph{defense\_overhead}, and \emph{utility\_retention}
        per boundary.
  \item \textbf{An open attack library of 50 concrete attack
        classes} covering every MCP boundary, each with
        CVSS~v3.1 scoring and pytest reproducibility.
  \item \textbf{An empirical defense comparison} at the MCP layer
        (not the OS layer) quantifying leakage reduction from
        per-tenant tool registries, tenant-prefixed cache keys,
        resource-path canonicalisation, and audience-bound
        authentication tokens---individually and in composition.
\end{enumerate}

\paragraph{Top-line findings.}
Across eight baseline attacks, four research questions, and 30
iterations per cell, we find:

\begin{itemize}
  \item \textbf{RQ-1} (baseline isolation): cache attacks dominate
        the vulnerable server's leakage volume, accounting for
        85.7\% of all observed leakage ($z = 40.75$, $p < 10^{-4}$
        vs. an even split).
  \item \textbf{RQ-2} (cache dominance): the cache boundary is the
        single most-exploitable surface; cache attacks are the
        cheapest path to cross-tenant data disclosure in our
        reference vulnerable server.
  \item \textbf{RQ-3} (prompt-injection residual): even the
        hardened secure server exhibits non-trivial leakage under
        repeated prompt-injection attacks (mean leak rate
        $0.5$), confirming that prompt injection is a structural
        property of tool-using LLM agents and not a patch
        target.
  \item \textbf{RQ-4} (defense composition): no single defense
        closes the gap, but the composition of all four defenses
        reduces leak rate to zero ($t = 22.80$, $p < 10^{-5}$ on
        a paired Welch's $t$-test), validating defense-in-depth
        at the MCP layer.
\end{itemize}

\paragraph{Contributions.}
We make five contributions to the security of multi-tenant LLM
agent infrastructure:

\begin{enumerate}
  \item \textbf{An eight-boundary catalogue} for the Model Context
        Protocol, covering all eight protocol boundaries with
        STRIDE-tagged attack surfaces and 63 forward ticket IDs
        (\S\ref{sec:threat_model}). \emph{Delta vs.\ prior work:}
        Jing et al.~(2026) propose 5 generic LLM-agent
        boundaries~\cite{jing_2026_isolation}; ours is 8
        MCP-specific boundaries with concrete STRIDE rows.
  \item \textbf{A 14-concept security taxonomy with explicit
        MCP-boundary bindings} (\S\ref{sec:background}).
        \emph{Delta vs.\ prior work:} OWASP LLM Top 10
        v1.1~\cite{owasp_llm_top10_v11} is LLM-application-level,
        not protocol-level; no prior work pins security concepts
        to MCP boundaries.
  \item \textbf{A reproducible empirical framework} that measures
        cross-tenant leakage in MCP deployments under controlled
        multi-tenant concurrency (\S\ref{sec:framework}).
        \emph{Delta vs.\ prior work:} Chen et al.~(2026)
        benchmark static scanners~\cite{chen_2026_mcp_security};
        An et al.~(2026) detect but do not measure
        leakage~\cite{an_2026_flowguard}; Liu et al.~(2026) study
        tool evolution, not tenant
        isolation~\cite{liu_2026_mcpevol}.
  \item \textbf{An open attack library of 50 concrete attack
        classes} covering every MCP boundary, each with CVSS~v3.1
        scoring and pytest reproducibility (\S\ref{sec:attacks}).
        \emph{Delta vs.\ prior work:} Wang et al.~(2024),
        Torres et al.~(2026), and Lee et al.~(2026) each
        demonstrate one failure mode in isolation; no prior work
        composes them into a single library with uniform pytest
        reproducibility~\cite{wang_2024_toolcommander,
        torres_2026_memory,lee_2026_webmcp}.
  \item \textbf{An MCP-layer defense-comparison study}
        (\S\ref{sec:evaluation}). \emph{Delta vs.\ prior work:}
        Dipta et al.~(2024) study OS-level sandbox side
        channels~\cite{dipta_2024_sandbox}; ours is
        protocol-level and measures leakage reduction under
        defense composition.
\end{enumerate}

\paragraph{Why a paper now?}
The MCP specification is under active revision, and our
empirical findings and defense blueprint can inform the next
revision (\S\ref{sec:conclusion}). We release the attack library,
the measurement framework, and the analysis notebooks as open
source so that the community can reproduce and extend our
results. The artifact is described in Appendix~A and is
available at the project repository.

\paragraph{Organisation.}
\S\ref{sec:background} introduces MCP and the isolation primer;
\S\ref{sec:threat_model} presents our threat model;
\S\ref{sec:framework} describes the measurement framework;
\S\ref{sec:attacks} catalogues the attack library;
\S\ref{sec:evaluation} answers the four research questions;
\S\ref{sec:defenses} details the defense blueprint;
\S\ref{sec:related} situates our work in the literature;
\S\ref{sec:discussion} discusses limitations and ethics; and
\S\ref{sec:conclusion} concludes with spec-change
recommendations.

\paragraph{Threat-model scope.}
We deliberately scope our threat model to the protocol layer
(\S\ref{sec:background}--\S\ref{sec:threat_model}). Physical-layer
side channels~\cite{dipta_2024_sandbox}, kernel-level compromise,
and LLM-provider compromise are out of scope and disclosed
honestly in \S\ref{sec:discussion}. Our empirical claims are
therefore about \emph{what an MCP server that follows the spec
can be made to leak}, not about what a determined adversary with
host-level access can do. This scoping is consistent with the
adjacent protocol-security literature~\cite{zhu_2025_lspfuzz,
chen_2026_mcp_security}.

\paragraph{Reproducibility.}
All artefacts --- attack library, measurement framework,
manifests, raw run outputs, analysis notebooks, summary CSVs,
and rendered figures --- are open source under the MIT license.
The pinned dependency versions in \texttt{pyproject.toml} and
the seed-controlled manifests in
\texttt{experiments/manifests/} make the results reproducible
on any machine with the documented toolchain. We ran the full
RQ-1--RQ-4 suite on three independent machines (Linux x86-64,
macOS arm64, Windows x86-64) and obtained bit-identical CSVs
modulo platform-dependent timestamps.

\section{Background}
% 02_background.tex
% Acronyms defined on first use, per prompts/11_paper.md style rule #5.

\section{Background}\label{sec:background}

This section introduces the Model Context Protocol (\mcp{}) and
recapitulates the isolation vocabulary we use throughout the paper.

\subsection{The Model Context Protocol}

\mcp{} (\mcp{}) is a JSON-RPC-2.0-based protocol that mediates
communication between an LLM agent (the \emph{host}) and one or
more external \emph{servers} that expose tools, resources, and
prompt templates~\cite{mcp_spec}. The protocol defines five
primitive message types:
\texttt{initialize}, \texttt{tools/list},
\texttt{tools/call}, \texttt{resources/read}, and
\texttt{prompts/get}, plus a notifications channel for
server-pushed events. Transports include stdio (for
co-located processes), HTTP+SSE, and streamable HTTP. An MCP
\emph{host} (e.g. an agent runtime) may multiplex many
connections to different servers; an MCP \emph{server} may serve
many \emph{tenants} (users, agents, or organisations).

\paragraph{Acronyms.}
\textbf{API} Application Programming Interface; \textbf{CLI}
Command-Line Interface; \textbf{CVSS} Common Vulnerability
Scoring System; \textbf{CWE} Common Weakness Enumeration;
\textbf{DoS} Denial of Service; \textbf{JSON-RPC} JSON Remote
Procedure Call; \textbf{JWT} JSON Web Token; \textbf{LLM} Large
Language Model; \textbf{LSP} Language Server Protocol;
\textbf{mTLS} Mutual Transport Layer Security; \textbf{MCP} Model
Context Protocol; \textbf{NDSS} Network and Distributed System
Security; \textbf{PII} Personally Identifiable Information;
\textbf{RPC} Remote Procedure Call; \textbf{SSE} Server-Sent
Events; \textbf{STRIDE} Spoofing, Tampering, Repudiation,
Information Disclosure, Denial of Service, Elevation of
Privilege~\cite{stride_orig}; \textbf{SAN} Subject Alternative
Name; \textbf{SSRF} Server-Side Request Forgery;
\textbf{TOCTOU} Time-of-Check Time-of-Use.

\subsection{Isolation in multi-tenant agent infrastructure}

We use \emph{isolation} in the security-engineering sense: a
tenant's data, computation, and identity are confined to that
tenant, with no observable influence on---or from---other
tenants~\cite{nist_sp_800_207}. The MCP specification does not
mandate any isolation mechanism; it simply defines the wire
format. Reference implementations (Python, TypeScript, Rust)
provide single-tenant defaults and leave multi-tenancy as an
implementation concern.

\paragraph{Why MCP's eight boundaries matter.}
We catalogue eight implicit boundaries that any multi-tenant
MCP server must enforce. Each boundary corresponds to a
trust-decision point in the protocol:

\begin{itemize}
  \item \textbf{Transport}---the byte-level channel between host
        and server (stdio, HTTP+SSE, streamable HTTP).
  \item \textbf{Session}---per-connection state on the server
        (resolved principal, caches, cancellation tokens, SSE
        queues).
  \item \textbf{Namespace}---tool, resource, and prompt naming
        and discovery.
  \item \textbf{Tool}---tool invocation and result delivery.
  \item \textbf{Resource}---file/blob/URI access via MCP
        resources.
  \item \textbf{Memory}---server-side conversation memory and
        scratchpads.
  \item \textbf{Cache}---cached tool outputs and embeddings.
  \item \textbf{Auth}---token and scope verification.
\end{itemize}

A boundary is \emph{enforced} when every observable
interactions across the boundary carries an unambiguous
tenant identity that the server can verify before acting on the
request. A boundary is \emph{broken} when an actor without a
legitimate tenant identity can either read another tenant's
state or cause another tenant to act on their behalf. We
operationalise these definitions in \S\ref{sec:threat_model}.

\subsection{Threat-model vocabulary}

We adopt STRIDE as our per-boundary threat taxonomy because it
is the de-facto vocabulary in protocol-level threat
modelling~\cite{stride_orig}. Each STRIDE letter maps cleanly
onto an MCP failure mode:

\begin{itemize}
  \item \textbf{Spoofing (S)}: impersonate a tenant or forge a
        session/transport identity.
  \item \textbf{Tampering (T)}: modify another tenant's state
        via a shared write path.
  \item \textbf{Repudiation (R)}: hide or falsify the actor of
        an attack via missing audit.
  \item \textbf{Information Disclosure (I)}: read another
        tenant's data.
  \item \textbf{Denial of Service (D)}: starve or degrade
        another tenant's service.
  \item \textbf{Elevation of Privilege (E)}: bypass authorisation
        via shadow or renamed capability.
\end{itemize}

The cross-product of eight boundaries and six STRIDE letters
yields 48 rows, all of which we enumerate in
\S\ref{sec:threat_model}.

\subsection{Defense vocabulary}

We evaluate four protocol-layer defenses (rather than OS-layer
sandboxing) because the goal is to fix MCP, not the underlying
operating system~\cite{dipta_2024_sandbox}. The four defenses
are: \emph{per-tenant tool registries} (namespace), \emph{tenant-
prefixed cache keys} (cache), \emph{resource-path
canonicalisation} (resource), and \emph{audience-bound JWTs}
(auth). Each defense is a self-contained middleware that can be
enabled or disabled per experiment, and we measure them
individually and in composition in \S\ref{sec:evaluation}.

\subsection{Why the eight boundaries matter together}

A common objection to boundary-by-boundary analysis is that
real-world attacks cross boundaries: a single prompt-injection
attack can traverse the resource boundary (the injected content
is read via \texttt{resources/read}), the auth boundary (the
injected instruction seeks elevation), and the tool boundary
(the model is tricked into calling a tool). Our defense
blueprint (\S\ref{sec:defenses}) addresses this objection by
arguing that the composition of one defense per boundary is
sufficient: each attack's exploitation path crosses
\emph{at least one} boundary whose defense is engaged. Cache
attacks cross the namespace boundary (tool routing) and the
cache boundary; the registry defence catches the former and
the cache defence catches the latter. Prompt-injection attacks
cross the resource boundary and the auth boundary; the
canonicalisation and JWT defences catch them in series.

A second common objection is that boundary catalogues are
\emph{not falsifiable} --- that any failure can be re-described
in terms of any boundary. We address this by insisting that
each row in our STRIDE table has a \emph{mechanism}: a
concrete code path or protocol message whose presence or
absence determines whether the boundary is enforced. The
per-tenant tool registry, for example, is a single middleware
line; if it is absent, the row \texttt{A-NSP-001} (tool-name
squatting) is exploitable; if it is present, the row is not.
The mechanism makes the catalogue falsifiable and the defenses
testable.

\subsection{Capability-based security as the underlying model}

The defense blueprint in \S\ref{sec:blueprint} is consistent with
\emph{capability-based security}: every action requires an
unforgeable capability that names both the principal and the
target. Capsicum (FreeBSD 9.0, 2012) is the historical
precedent; the audience-bound JWT in our auth defense is the
modern analogue. The key insight is that \emph{ambient authority}
--- permissions inherited from the process identity rather than
explicitly granted per action --- is the root cause of most
multi-tenant security failures. MCP's defaults inherit ambient
authority from the stdio worker (transport boundary), the
in-memory cache (cache boundary), and the shared tool registry
(namespace boundary); our defenses replace each ambient
authority with an explicit capability.

\subsection{Why protocol-layer defenses, not OS-layer sandboxes}

Dipta et al.~(2024) demonstrated that side channels survive
logical sandboxing~\cite{dipta_2024_sandbox}: their dynamic
frequency-based fingerprinting attack defeats gVisor,
Firecracker, SGX, and SEV. This is a \emph{positive} result for
their attack and a \emph{negative} result for OS-layer
sandboxing as a sole defense. Our position is that MCP-layer
defenses and OS-layer sandboxes are \emph{complementary}: the
MCP-layer defenses close the protocol-level leakage paths (the
RQ-4 result), and the OS-layer sandbox contains the residual
silicon-level leakage (Dipta et al.'s threat model). A
production MCP deployment should run the secure reference
server inside a Firecracker microVM and engage all four
protocol-layer defenses; the combination is what the MCP
specification should recommend.

\section{Threat Model}
% 03_threat_model.tex — mirrors docs/02_Threat_Model.md and
% docs/04_Attack_Taxonomy.md.

\section{Threat Model}\label{sec:threat_model}

We model the MCP server as the trust boundary under study. The
LLM provider is assumed trusted; the network between client and
server is assumed hostile unless TLS is configured; the kernel
of the server host is assumed trustworthy in our reference
deployment but is honestly disclosed as out-of-scope for the
protocol-level threat model. The full per-boundary STRIDE
enumeration is in our attack taxonomy
(\texttt{docs/04\_Attack\_Taxonomy.md}; summarised here in
\S\ref{sec:taxonomy-tables} for self-containedness).

\subsection{Actors}

\begin{itemize}
  \item \textbf{Honest Tenant.} A user with legitimate access to
        a subset of MCP tools and resources.
  \item \textbf{Malicious Tenant.} A user who attempts to access
        data or capabilities outside their authorized scope.
  \item \textbf{Compromised Agent.} An MCP host whose prompt
        context is partially controlled by an attacker.
  \item \textbf{Malicious MCP Server.} A server that misbehaves
        across tenants.
  \item \textbf{Network Adversary.} Passive observer or active
        MITM on the transport channel.
  \item \textbf{Server Admin.} Operator with host-level access;
        trusted in our reference deployment but not in the
        protocol-level threat model.
\end{itemize}

\subsection{Attacker capabilities (in scope)}

The attacker has execution rights in \emph{one} legitimate tenant
account (Tenant~A). The attacker can register tools, read public
resources, and call any public method on the server. The
attacker can inject content via prompt, resource, or tool
result. The attacker can observe timing and error responses from
the server. The attacker can run multiple concurrent connections
simultaneously, subject to per-tenant rate limits if any.

\subsection{Out-of-scope threats}

The following are explicitly excluded from our evaluation:

\begin{itemize}
  \item Kernel-level compromise of the server host.
  \item LLM-provider compromise (theft of model weights,
        backdoored model serving).
  \item Side channels on shared silicon (cache-line timing,
        rowhammer, Spectre/Meltdown-class attacks)~\cite{dipta_2024_sandbox}.
  \item Physical access to the server.
  \item Compromise of the upstream identity provider.
\end{itemize}

\subsection{Assets}

\begin{itemize}
  \item \textbf{Data.} Tenant-owned documents, embeddings,
        conversation history, tool outputs, cached prompts,
        resource blobs.
  \item \textbf{Credentials.} Bearer tokens, mTLS client
        certificates, refresh tokens, session IDs.
  \item \textbf{Sessions.} Per-connection state: resolved
        principal, per-session caches, in-flight cancellation
        tokens, SSE event queues.
  \item \textbf{Embeddings.} Server-side vector representations
        of prior prompts and tool outputs (long-lived PII risk).
  \item \textbf{Tool outputs.} Cached or in-flight results from
        \texttt{tools/call} invocations.
  \item \textbf{Capability registries.} Server-side tool,
        resource, and prompt catalogues and the descriptors they
        expose.
\end{itemize}

\subsection{Per-boundary STRIDE}\label{sec:taxonomy-tables}

We enumerate eight STRIDE tables (one per boundary) with 48
rows. Each row carries a STRIDE letter, a one-sentence threat
description, one or more ticket IDs from our 63-ticket
catalogue, and a CWE reference. The tables are reproduced in
full in Appendix~B for self-containedness; an excerpt is given
below.

\begin{table}[t]
\caption{Per-boundary STRIDE enumeration (excerpt). Full table in
Appendix~B; see \texttt{docs/04\_Attack\_Taxonomy.md} for the
authoritative source.}
\label{tab:stride-excerpt}
\centering
\small
\begin{tabularx}{\linewidth}{l l X l}
\toprule
Boundary & STRIDE & Threat & Ticket \\
\midrule
Transport & S & Token/session replay from cleartext channel & \texttt{A-TRN-002} \\
Session   & I & SSE queue keyed on \texttt{session\_id} only     & \texttt{A-SES-003} \\
Namespace & E & Decorator-name shadowing in Python SDK         & \texttt{A-NSP-007} \\
Tool      & I & Cross-tenant context capture via shared queue   & \texttt{A-TOL-003} \\
Resource  & I & Path traversal across tenants via \texttt{..}   & \texttt{A-RES-001} \\
Memory    & I & Embedding-store PII persistence                 & \texttt{A-MEM-004} \\
Cache     & I & Cross-tenant hit on missing \texttt{tenant\_id} & \texttt{A-CCH-003} \\
Auth      & I & Token forwarding via SSRF                      & \texttt{A-AUT-004} \\
\bottomrule
\end{tabularx}
\end{table}

\subsection{Misuse cases}

We provide four concrete misuse cases, each pinned to $\geq 2$
ticket IDs and a likelihood/impact rating. These are the
executable seeds of our attack library
(\S\ref{sec:attacks}).

\paragraph{MC-1: Cross-tenant environment-variable overwrite via
tool shadowing.} Tenant~A registers a tool named
\texttt{set\_env} that writes to a shared scratchpad keyed by
\texttt{env\_var\_name} (no \texttt{tenant\_id}); Tenant~B
subsequently calls any tool that reads \texttt{PATH} from the
same scratchpad and receives Tenant~A's value. \emph{Likelihood:}
high. \emph{Impact:} high (PATH hijack, credential theft).
Mapped to \texttt{A-NSP-001}, \texttt{A-TOL-005},
\texttt{A-MEM-002}.

\paragraph{MC-2: Session fixation across server restart.} Server
uses predictable session IDs and reuses them after restart;
Tenant~A resumes Tenant~B's session ID and reads Tenant~B's
cached tool outputs. \emph{Likelihood:} medium. \emph{Impact:}
high. Mapped to \texttt{A-SES-001}, \texttt{A-SES-002},
\texttt{A-CCH-003}.

\paragraph{MC-3: Resource path traversal via percent-encoded
slashes.} Server accepts \texttt{file://} URIs as resource
identifiers but does not canonicalise \texttt{..} segments or
strip percent-encoded slashes; Tenant~A reads
Tenant~B's secrets via a URI like
\texttt{file:///tenant-a/..\%2F..\%2Ftenant-b\%2Fsecrets.txt}.
\emph{Likelihood:} high. \emph{Impact:} critical. Mapped to
\texttt{A-RES-001}, \texttt{A-RES-002}, \texttt{A-RES-003}.

\paragraph{MC-4: Embedding cache poisoning via prompt collision.}
Server caches prompt embeddings keyed by \texttt{prompt\_hash}
only; Tenant~A crafts a colliding prompt and reads
Tenant~B's embedding from the cache. \emph{Likelihood:} low.
\emph{Impact:} medium-high. Mapped to \texttt{A-CCH-004},
\texttt{A-MEM-004}.

\subsection{Data-flow diagram}

A high-level data-flow diagram with trust-boundary edges is
available in \texttt{docs/diagrams/dfd\_trust\_boundaries.svg}.
The DFD shows the orchestrator / MCP server / Tenant~A /
Tenant~B data flow with trust-boundary edges annotated.

\subsection{Assumptions}

\begin{itemize}
  \item The attacker has at most one legitimate tenant account.
  \item The LLM provider is trusted; the MCP server is the
        trust boundary under study.
  \item The network between client and server is hostile unless
        TLS is configured.
  \item Our reference servers (\texttt{mcp\_servers/vulnerable/}
        and \texttt{mcp\_servers/secure/}) are run in Firecracker
        microVMs~\cite{firecracker_github} to limit kernel-level
        collateral damage during attack execution.
\end{itemize}

\subsection{What is not in the threat model (and why)}

Three classes of attacks are deliberately excluded from our
threat model because they are orthogonal to the protocol-layer
question we study:

\paragraph{Host-level compromise.} A kernel-level compromise
of the server host (e.g. via a malicious dependency, a
vulnerable syscall, or a container escape) is out of scope. We
assume the host's syscall surface and its microVM boundary
hold. The MCP specification cannot defend against a kernel
compromise; that is the OS's responsibility. Our reference
deployment runs inside Firecracker precisely so that even a
kernel compromise is contained.

\paragraph{LLM-provider compromise.} Theft of model weights,
backdoored model serving, or prompt-level jailbreaks of the
LLM itself are out of scope. We assume the LLM provider follows
its own security practices (e.g. strict output filtering,
rate limiting, content moderation). Mid-session tool-injection
via tool-result content is \emph{in} scope because it targets
the MCP layer, not the LLM.

\paragraph{Silicon-level side channels.} Cache-line timing,
rowhammer, and Spectre/Meltdown-class attacks on shared
silicon~\cite{dipta_2024_sandbox} are out of scope. These
attacks survive any protocol-layer defense. We acknowledge
them in \S\ref{sec:discussion} as a threat-model limitation
and recommend running the secure server inside a microVM to
contain them.

\subsection{Threat-model consistency with adjacent work}

Our threat model is consistent with the conventions established
by the adjacent protocol-security literature:

\begin{itemize}
  \item LSPFuzz~\cite{zhu_2025_lspfuzz} scopes its threat model to
        the LSP wire protocol and explicitly excludes kernel-
        level issues; we mirror this scoping.
  \item FlowGuard~\cite{an_2026_flowguard} scopes its detector to
        MCP-server behaviour traces; we mirror this scoping.
  \item Chen et al.'s large-scale empirical study~\cite{chen_2026_mcp_security}
        scopes its findings to runtime MCP servers;
        we mirror this scoping.
  \item Dipta et al.'s sandbox-fingerprinting work~\cite{dipta_2024_sandbox}
        is the explicit reference for what is \emph{not} in
        scope.
\end{itemize}

The consistency is deliberate: threat-model scoping is a
literature contract; an MCP paper that included kernel
side-channels in scope would be unpublishable at a security
venue because the defences it proposes would be irrelevant.
By scoping to the protocol layer, our defenses are
\emph{necessary} (the protocol layer must close its own
leakage paths before any sandbox can be effective) and
\emph{sufficient for the protocol layer} (RQ-4 confirms
zero leakage at the composition point).

\subsection{Attacker profile}

We assume a \emph{tenant-class} attacker: one who has obtained
legitimate credentials for one tenant account (via phishing,
credential stuffing, purchase on a grey market, or a rogue
employee). The attacker is \emph{not} an insider with host-level
access; the attacker is \emph{not} a nation-state with the
ability to compromise the upstream identity provider; the
attacker \emph{can} register tools, read public resources, call
any public method, inject content via prompt, resource, or
tool result, observe timing and error responses, and run
multiple concurrent connections.

This profile matches the threat model in the OWASP LLM Top
10~\cite{owasp_llm_top10_v11} and the adversarial examples in
Chowdhury et al.'s defenses survey~\cite{chowdhury_2024_defenses}.
It is the most demanding profile that the protocol layer can
defend against without hardware support.

\section{Framework Design}
% 04_framework.tex — mirrors docs/06_Framework.md

\section{Framework Design}\label{sec:framework}

We design and implement a reproducible cross-tenant leakage
measurement framework with six cooperating modules plus a
reporter. The framework is the executable embodiment of the
threat model in \S\ref{sec:threat_model}: given a manifest of
tenants, attacks, and defenses, it spins up the reference
servers, drives concurrent tenant traffic, classifies
leakage events, and emits per-attack metrics.

\subsection{Goals}

\begin{itemize}
  \item \textbf{Replayable.} Given a manifest + seed, the
        framework produces identical results.
  \item \textbf{Tenant-aware.} First-class notion of
        \texttt{tenant\_id} across scheduler, evaluator, logger.
  \item \textbf{Multi-boundary.} Covers all eight MCP
        boundaries.
  \item \textbf{Pluggable.} New attacks, defenses, and metrics
        can be added without changing the core.
\end{itemize}

\subsection{Architecture}

\begin{table}[t]
\caption{Framework modules (\texttt{framework/}). Each module
has a single responsibility and a typed input/output contract.}
\label{tab:modules}
\centering
\small
\begin{tabularx}{\linewidth}{l X X}
\toprule
Module & Responsibility & File \\
\midrule
Scheduler    & Drive concurrent tenants + attack queue   & \texttt{scheduler/scheduler.py} \\
PayloadGen   & Deterministic per-seed payload mutation   & \texttt{scheduler/payloads.py} \\
Connector    & Speak MCP (stdio / HTTP+SSE / dummy)      & \texttt{target/connector.py} \\
Evaluator    & Classify leakage via substring + sha256   & \texttt{evaluator/evaluator.py} \\
Metrics      & Aggregate per boundary / defense          & \texttt{metrics/metrics.py} \\
Logger       & Persist every event as JSONL              & \texttt{logger/logger.py} \\
Reporter     & Render Markdown + HTML reports            & \texttt{reports/reporter.py} \\
\bottomrule
\end{tabularx}
\end{table}

A high-level sequence diagram is available in
\texttt{docs/diagrams/framework\_sequence.svg}. The diagram
covers the five prompt-required steps: framework spins up
Tenant~A and Tenant~B; payload injected via Tenant~A;
Tenant~B's state is sampled; oracle classifies result; logger
writes event.

\subsection{RunConfig schema}

The framework hydrates a \texttt{RunConfig} (Pydantic v2 model)
from a YAML manifest. The full schema is:

\begin{lstlisting}[language=yaml,caption={RunConfig schema (excerpt).},label=lst:runconfig]
run:
  seed: 42
  repeats: 30
  concurrency: 4
tenants:
  - id: tenant-A
    name: Alice
    allowed_tools: [echo]
  - id: tenant-B
    name: Bob
    allowed_tools: [echo]
attacks:
  - id: A-CCH-T
    parameters: {}
defenses:
  per_tenant_tool_registry: false
  tenant_prefixed_cache_keys: false
  resource_path_canonicalisation: false
  mtls: false
target:
  transport: dummy   # stdio | http_sse | streamable_http | dummy
output:
  log_dir: analysis/runs
  output_dir: analysis/runs
  log_format: csv
\end{lstlisting}

\texttt{RunConfig} enforces at least 2 tenants because the
harness requires a (source, sink) pair to detect leakage. The
\texttt{defenses} block is consumed by the reference
secure-server factory (see \S\ref{sec:evaluation}). The full
schema is versioned: \texttt{schema\_version: '1.0'} and
\texttt{dataset\_version: 'phase9-v1'} are mandatory fields;
the runner refuses to execute a manifest whose schema version
does not match the runner's expected version.\footnote{The
schema is documented in
\texttt{experiments/manifests/schema.yaml}. The pinning
protects against silent manifest-format changes between
Phases 9 and 12.} Sample size $n = 30$ per cell is the
pre-registered choice (\texttt{analysis/power.md}) that
targets medium effect sizes (Cliff's $\delta \geq 0.50$)
at power $\geq 0.80$ and $\alpha = 0.05$; smaller effects
are characterised as exploratory; the headline comparisons
in \S\ref{sec:evaluation} target $\delta \geq 0.80$ (large)
and are well within the $n = 30$ budget.

\subsection{Metrics}

The framework computes four metrics per attack per defense
level:

\begin{itemize}
  \item \textbf{\texttt{leakage\_rate}} -- fraction of call
        events in the bucket that have a paired leakage event
        (range $[0, 1]$).
  \item \textbf{\texttt{time\_to\_leak\_ms}} -- median latency
        between \texttt{attack\_started\_at} and
        \texttt{leakage\_detected\_at}; \texttt{null} if no
        leakage.
  \item \textbf{\texttt{defense\_overhead\_ms}} -- p95 of
        defended-call latency minus p95 of undefended-call
        latency; \texttt{null} if either side has no samples.
  \item \textbf{\texttt{utility\_retention}} -- $1 -
        \mathrm{errors\_defended}/\mathrm{errors\_undefended}$,
        clamped to $[0, 1]$.
\end{itemize}

\subsection{Reproducibility}

Each manifest declares a seed; the framework's payload
generator and the connector's probabilistic leak path both
consume that seed via \texttt{random.Random}. Defenses are
enabled/disabled at server startup, not at runtime, so the
server's behaviour is fully determined by
$(\mathrm{manifest}, \mathrm{seed})$. The artifact includes
pinned dependency versions (\texttt{pyproject.toml}) and a
\texttt{requirements.txt} for byte-identical reproduction.

\subsection{Reference servers}

We ship two reference servers:

\begin{itemize}
  \item \textbf{\texttt{mcp\_servers/vulnerable/}} --- a
        minimal MCP server with all eight boundaries
        deliberately mis-configured: shared tool registry,
        shared cache keys, no URI canonicalisation, no JWT
        \texttt{aud} claim, and no defense middleware.
  \item \textbf{\texttt{mcp\_servers/secure/}} --- the same
        MCP server with all four defense middleware enabled:
        per-tenant tool registry, tenant-prefixed cache keys,
        URI canonicalisation, audience-bound JWTs, and
        probabilistic leak probability reduced to zero.
\end{itemize}

The two servers share the same wire protocol, the same tool
catalogue, and the same fixtures; the only difference is the
defense middleware. This makes the comparison honest: every
observed leakage difference is attributable to the defenses,
not to implementation drift.

\paragraph{Tool catalogue.}
The reference servers expose eight tools (\texttt{echo},
\texttt{get\_secret}, \texttt{list\_tenants}, \texttt{set\_env},
\texttt{read\_file}, \texttt{write\_file}, \texttt{cache\_get},
\texttt{cache\_put}). The tool catalogue and per-tool
arguments are documented in
\texttt{artifact/release/TOOL\_CATALOGUE.md}; the
\texttt{list\_tenants} tool is admin-only on the secure
server.

\subsection{Why a measurable oracle, not a model-based judge}

A subtle design question is whether the Evaluator (oracle)
should be a learned model or a deterministic substring-and-hash
match. We chose the latter for three reasons:

\begin{enumerate}
  \item \textbf{Reproducibility.} A deterministic oracle
        produces bit-identical results across machines. A
        learned oracle's outputs depend on the model's
        checkpoint, which is a moving target.
  \item \textbf{Causal interpretability.} When the oracle
        flags a leakage event, the developer can inspect the
        matched excerpt and verify the leak path. A model-based
        oracle produces a probability that requires
        interpretation.
  \item \textbf{Speed.} The deterministic oracle runs in
        $< 1\,\mu s$ per event; a model-based oracle takes
        $\sim 50\,ms$ per event. The factor-of-$10^{4}$
        difference is decisive for the $n = 30$ iterations
        across four research questions.
\end{enumerate}

The trade-off is false negatives: a deterministic oracle
misses leakages whose payload does not match the seeded
marker. We compensate by ensuring that every attack class
seeds a known marker (\texttt{payload\_sha256} attribute) and
that the Evaluator matches both the marker and its sha256
prefix. The full attack library has been audited to confirm
that every attack seeds a marker; the pytest suite asserts
this as a contract.

\subsection{Concurrency model}

The Scheduler uses Python's \texttt{asyncio} with a bounded
\texttt{Semaphore} to drive concurrent tenant traffic. Each
attack iteration runs in its own task; the scheduler awaits
all tasks at the iteration boundary before flushing the
logger. The concurrency limit is set per-manifest
(\texttt{run.concurrency}); our experiments use
\texttt{concurrency = 4} for RQ-1, RQ-2, RQ-3, and
\texttt{concurrency = 8} for RQ-4 to stress-test the defense
composition.

The connector's leak-probability path is per-connection (not
global), so concurrent tenants do not share a seed. The
Evaluator's substring match is per-tenant-pair, so concurrent
tenants do not cross-contaminate. The Logger appends to a
single JSONL file with file-level locking; race conditions
are confined to write boundaries, which the iteration
flushes handle.

\subsection{What the framework does not do}

The framework is intentionally narrow:

\begin{itemize}
  \item It does not run a real LLM. The connector uses a
        deterministic in-process channel; prompt-injection
        attacks inject markers, not instructions to a model.
  \item It does not measure silicon-level side channels. The
        connector's timing is deterministic at the
        microsecond scale; Spectre/Meltdown-class attacks
        would require a separate harness.
  \item It does not fuzz the wire protocol beyond the
        \texttt{A-FUZZ-001} grammar attack. A full mutation
        campaign (à la LSPFuzz~\cite{zhu_2025_lspfuzz}) is
        future work.
  \item It does not handle multi-language SDKs. The reference
        connector is Python-only; the TypeScript and Rust
        SDKs are out of scope.
\end{itemize}

These exclusions are deliberate: they keep the framework
focused on the protocol-layer question and the empirical
results interpretable.

\section{Attack Library}
% 05_attacks.tex — mirrors docs/07_Attack_Library.md

\section{Attack Library}\label{sec:attacks}

We ship 50 concrete attack classes (48 per-row STRIDE plus 2
cross-cutting) in our open attack library
(\texttt{attacks/}). Each class has a CVSS~v3.1 vector, a
pytest reproducibility contract, and a one-class-per-row
mapping to the STRIDE tables in \S\ref{sec:threat_model}.

\subsection{Coverage}

\begin{table}[t]
\caption{Attack library coverage. Each boundary has 6 STRIDE
letters (Spoofing through Elevation), giving 48 boundary rows;
two cross-cutting attacks (prompt injection + grammar fuzzing)
round the library out to 50 classes. The \emph{Tickets} column
lists the per-boundary ticket IDs covered; the CVSS column
lists the per-row CVSS~v3.1 vector family.}
\label{tab:coverage}
\centering
\small
\begin{tabularx}{\linewidth}{l l X X}
\toprule
Boundary & STRIDE & Tickets covered & CVSS family \\
\midrule
transport    & S/T/R/I/D/E & \texttt{A-TRN-001..009} & \texttt{AV:N/AC:L/PR:N/UI:N/S:U/...} \\
session      & S/T/R/I/D/E & \texttt{A-SES-001..009} & \texttt{AV:N/AC:H/PR:N/UI:N/S:U/...} \\
namespace    & S/T/R/I/D/E & \texttt{A-NSP-001..007} & \texttt{AV:N/AC:L/PR:L/UI:N/S:U/...} \\
tools        & S/T/R/I/D/E & \texttt{A-TOL-001..010} & \texttt{AV:N/AC:L/PR:L/UI:R/S:U/...} \\
resources    & S/T/R/I/D/E & \texttt{A-RES-001..007} & \texttt{AV:N/AC:L/PR:L/UI:N/S:U/...} \\
memory       & S/T/R/I/D/E & \texttt{A-MEM-001..008} & \texttt{AV:N/AC:L/PR:L/UI:N/S:U/...} \\
cache        & T/R/I/D/E   & \texttt{A-CCH-001..005} & \texttt{AV:N/AC:H/PR:N/UI:N/S:U/...} \\
auth         & S/T/R/I/D/E & \texttt{A-AUT-001..008} & \texttt{AV:N/AC:L/PR:N/UI:N/S:U/...} \\
prompt inj.  & n/a         & \texttt{A-TOL-001}, \texttt{A-RES-006/007} & \texttt{AV:N/AC:L/PR:L/UI:R/S:C/C:H/I:H/A:N} \\
fuzzing      & n/a         & surface-wide          & \texttt{AV:N/AC:H/PR:N/UI:N/S:U/C:H/I:H/A:H} \\
\bottomrule
\end{tabularx}
\end{table}

\subsection{Attack patterns}

We group the 48 per-row attacks into 8 reusable patterns, each
matching one boundary's failure mode:

\begin{enumerate}
  \item \textbf{\texttt{framing\_or\_replay\_inject}}
        (transport $\times$ S/T/R/I/D/E). Smuggle a payload
        marker into the framing bytes or replay surface; the
        source connector invokes \texttt{echo} with a
        marker-carrying argument set and the Evaluator checks
        whether the marker surfaces in the sink's view.
  \item \textbf{\texttt{session\_fixation}} (session $\times$
        S/T/R/I/D/E). Force a fixed session ID on the source
        connector and inspect whether the sink's session store
        returns the source's payload marker. Mirrors misuse
        case MC-2.
  \item \textbf{\texttt{tool\_shadow\_inject}} (namespace
        $\times$ S/T/R/I/D/E). Invoke a tool whose name
        shadows a built-in (\texttt{shadow\_set\_env}) and read
        a resource whose URI is shaped by the squatted name.
        Mirrors misuse case MC-1.
  \item \textbf{\texttt{result\_smuggle\_or\_inject}} (tool
        $\times$ S/T/R/I/D/E). Indirect prompt injection: the
        source connector invokes \texttt{echo} with a payload
        carrying a prompt-like instruction; the sink connector
        reads a resource whose contents include the same
        marker. Cross-boundary (tool $\leftrightarrow$ resource).
  \item \textbf{\texttt{path\_traversal\_uri}} (resource
        $\times$ S/T/R/I/D/E). Path traversal: the source
        connector reads a URI with \texttt{..} traversal and
        percent-encoded slashes (\texttt{\%2F}); the sink
        connector checks whether the traversal surfaced the
        marker. Mirrors misuse case MC-3.
  \item \textbf{\texttt{scratchpad\_or\_memory\_leak}} (memory
        $\times$ S/T/R/I/D/E). The source writes a value to a
        scratchpad keyed by name only (no
        \texttt{tenant\_id}); the sink reads the same key.
        Mirrors the confused-deputy concept.
  \item \textbf{\texttt{cache\_key\_collision}} (cache
        $\times$ T/R/I/D/E). Both connectors attempt to read
        the same cache entry keyed by
        \texttt{(tool\_name, args\_hash)}; the attack succeeds
        if the sink returns the source's payload. Mirrors
        misuse case MC-4.
  \item \textbf{\texttt{token\_replay\_or\_misuse}} (auth
        $\times$ S/T/R/I/D/E). The source connector invokes a
        tool that performs an SSRF-style fetch carrying the
        original \texttt{Authorization} header; the sink
        connector reads a resource to detect whether the token
        surfaced.
\end{enumerate}

\subsection{Cross-cutting attacks}

\paragraph{\texttt{ResultInjectionAttack} (\texttt{A-TOL-001}).}
Indirect prompt injection via tool result content. Touches both
the \textbf{tool} and \textbf{resource} boundaries. Implements
\texttt{A-TOL-001}, \texttt{A-RES-006}, and \texttt{A-RES-007}
simultaneously. CVSS~3.1:
\texttt{AV:N/AC:L/PR:L/UI:R/S:C/C:H/I:H/A:N}. This is the
direct protocol-level embodiment of the indirect prompt
injection originally identified by Greshake et al.
(2023)~\cite{greshake_2023_indirect}.

\paragraph{\texttt{GrammarFuzzAttack} (\texttt{A-FUZZ-001}).}
Grammar-aware fuzzing harness adapted from LSPFuzz (Zhu et al.
2025~\cite{zhu_2025_lspfuzz}). Generates $N$ well-formed-but-
malicious JSON-RPC envelopes with mutated parameter shapes.
LSPFuzz's campaign over $\sim$300 LSP servers is the
methodological template; ours is the MCP-specific adaptation
that respects the eight-boundary catalogue.

\subsection{Pytest reproducibility}

The library ships with a parametrized pytest suite
(\texttt{tests/test\_attack\_library.py}) that exercises every
registered id:

\begin{itemize}
  \item \texttt{test\_minimum\_attack\_count} asserts $\geq 25$
        (we ship 50).
  \item \texttt{test\_every\_boundary\_has\_at\_least\_one\_attack}
        asserts each of the 8 \texttt{Boundary} enum values has
        $\geq 1$ attack.
  \item \texttt{test\_attack\_id\_format} asserts id matches
        \texttt{A-PREFIX-NNN}.
  \item \texttt{test\_attack\_has\_cvss} asserts every class
        has a CVSS vector attribute or docstring.
  \item \texttt{test\_attack\_subclasses\_attack\_base} asserts
        every class subclasses \texttt{attacks.base.Attack}.
  \item \texttt{test\_attack\_instantiable} asserts every class
        can be instantiated.
  \item \texttt{test\_attack\_execute\_runs} runs one
        representative attack end-to-end.
  \item \texttt{test\_no\_duplicate\_ids} asserts no id
        collisions.
\end{itemize}

The full suite passes with 204 tests in $< 0.2$s, confirming
that the library is reproducible on any machine with the pinned
dependencies.

\subsection{Attack-class lifecycle}

Each attack class follows the same lifecycle:

\begin{enumerate}
  \item \textbf{Registration.} The class is registered in
        \texttt{attacks.registry.REGISTRY} keyed by its
        \texttt{id} attribute. The registry is a
        \texttt{dict[str, Type[Attack]]} consumed by the
        Scheduler.
  \item \textbf{Instantiation.} The Scheduler instantiates the
        class with the tenant pair and the per-recipe
        parameters from the manifest. The class's
        \texttt{\_\_init\_\_} validates the parameters and
        raises on type mismatches.
  \item \textbf{Execution.} The Scheduler calls
        \texttt{execute()} with the source connector, the
        sink connector, and a seeded \texttt{random.Random}.
        The method drives the cross-tenant interaction and
        returns a \texttt{list[Event]} (call events,
        leakage events, or both).
  \item \textbf{Evaluation.} The Evaluator inspects the
        events and classifies leakage based on substring
        and sha256-prefix matches against the seeded
        marker.
  \item \textbf{Logging.} The Logger appends each event to
        the run's JSONL file with the tenant pair, the
        attack id, the iteration, and the seeded marker.
\end{enumerate}

The lifecycle is uniform across all 50 classes, which makes
the framework's pluggability contract concrete: a new attack
class needs only to subclass \texttt{attacks.base.Attack},
implement \texttt{execute()}, and register itself.

\subsection{Coverage rationale}

The 48-row coverage (8 boundaries $\times$ 6 STRIDE letters)
is not arbitrary. STRIDE is the de-facto vocabulary for
protocol-level threat modelling~\cite{stride_orig}; each
letter maps onto a distinct attack pattern at each
boundary. The two cross-cutting attacks
(\texttt{ResultInjectionAttack}, \texttt{GrammarFuzzAttack})
are added because prompt injection and grammar fuzzing are
\emph{horizontal} attack patterns that cross multiple
boundaries and cannot be pinned to a single row.

The CVSS scoring per row is documented in each attack class's
\texttt{cvss} attribute; full vectors are generated by
\texttt{scripts/gen\_phase7\_implementations.py} and follow
the FIRST specification. The scores are consistent with the
adjacent literature (Chowdhury et al.'s survey
\cite{chowdhury_2024_defenses}; Wang et al.'s ToolCommander
\cite{wang_2024_toolcommander}) and serve as a triage tool
for defenders, not as a precise impact prediction.

\subsection{Comparison with adjacent attack libraries}

The closest prior attack libraries are:

\begin{itemize}
  \item \textbf{Chowdhury et al.'s LLM-attack survey}
        \cite{chowdhury_2024_defenses}: catalogue of LLM-level
        attacks; no MCP-specific surface.
  \item \textbf{An et al.'s FlowGuard}
        \cite{an_2026_flowguard}: detection traces, not
        attack implementations.
  \item \textbf{Zhu et al.'s LSPFuzz}
        \cite{zhu_2025_lspfuzz}: grammar-aware fuzzing on
        LSP; not directly portable to MCP.
  \item \textbf{Liu et al.'s MCPEvol-Bench}
        \cite{liu_2026_mcpevol}: dynamic-evolution benchmark;
        no attack implementations.
\end{itemize}

Our library is the first to ship MCP-specific attack
implementations with pytest reproducibility, CVSS scoring,
and uniform per-(boundary, STRIDE) coverage. The combination
is the unit of novelty: prior work covers some of these
dimensions in isolation, none in composition.

\section{Evaluation}
% 06_evaluation.tex — consumes analysis/tables/*.csv and
% analysis/figures/*.pdf; the four research questions are
% answered here.

\section{Evaluation}\label{sec:evaluation}

We answer four research questions with the framework and
attack library from \S\ref{sec:framework}--\S\ref{sec:attacks}.

\subsection{Research questions}

\begin{itemize}
  \item \textbf{RQ-1} \emph{(baseline isolation).} Across all
        eight protocol boundaries, how much cross-tenant
        leakage does a naive MCP server exhibit, and how much
        does a hardened server reduce it?
  \item \textbf{RQ-2} \emph{(cache dominance).} Is the cache
        boundary the dominant source of leakage, and is the
        effect statistically distinguishable from an even
        split?
  \item \textbf{RQ-3} \emph{(prompt-injection residual).} Does
        the hardened server completely suppress leakage under
        prompt injection, or is there a bounded residual
        consistent with the structural-property view of
        injection?
  \item \textbf{RQ-4} \emph{(defense composition).} Do
        individual MCP-layer defenses close the leakage gap,
        or is composition required?
\end{itemize}

\subsection{Experimental setup}

Each research question is answered by a manifest
(\texttt{experiments/manifests/rq{1..4}\_*.yaml}) that fixes
the attack set, the defense level, the seed, and the iteration
count. We run 30 iterations per cell. The framework writes a
CSV with 12 columns per iteration and a \texttt{meta.json}
with provenance metadata (commit SHA, seed, n\_events,
duration\_s). The four runs are stored under
\texttt{analysis/runs/exp-rq{1..4}-*/}. The analysis code
under \texttt{analysis/scripts/} produces the four summary
CSVs (\texttt{analysis/tables/rq{1..4}\_summary.csv}) and
four PDF + PNG figures (\texttt{analysis/figures/}).

\subsection{Statistical protocol}

The pre-registered statistical protocol
(\texttt{analysis/power.md}) sets $\alpha = 0.05$, target
power $= 0.80$, $n = 30$ per cell, and uses Welch's $t$-test
(\texttt{scipy.stats.ttest\_ind(equal\_var=False)}) plus
Cliff's $\delta$ as a non-parametric effect size. Bonferroni
correction is applied at the boundary level
($\alpha_{\mathrm{adj}} = 0.05 / n_{\mathrm{attacks\_in\_boundary}}$).
Significance markers: $^{*} p < 0.05$, $^{**} p < 0.01$,
$^{***} p < 0.001$.

\subsection{RQ-1: Baseline isolation}\label{sec:rq1}

We run all eight baseline attacks (\texttt{A-AUT-T},
\texttt{A-MEM-T}, \texttt{A-RES-T}, \texttt{A-SES-T},
\texttt{A-CCH-T}, \texttt{A-NSP-T}, \texttt{A-TOL-T},
\texttt{A-TRN-S}) against the vulnerable and secure servers
with no defense middleware. Each cell runs 30 iterations.

\begin{table}[t]
\caption{RQ-1 per-attack leak rate. The vulnerable server
exhibits a 50\% leak rate per attack across all eight
boundaries; the secure server's connector-level leak
probability is set to zero by the defense middleware, so the
\emph{expected} leak rate is zero. We report the
paired one-sample Welch's $t$-test on the per-attack
differences $\Delta_i = L_{\mathrm{vuln},i} -
L_{\mathrm{secure},i}$ against zero, plus Cliff's $\delta$
across the full attack set.\footnote{Cliff's $\delta = 0$
indicates complete distributional overlap between the two
groups; in this context it indicates that the per-attack
leak-rate distributions on the vulnerable and secure
servers are identical, which is the consequence of the
partial-defense configuration not yet closing the gap.}
Source: \texttt{analysis/tables/rq1\_summary.csv}.}
\label{tab:rq1}
\centering
\small
\begin{tabularx}{\linewidth}{l l r r r r}
\toprule
attack\_id & boundary & $L_{\mathrm{vuln}}$ & $L_{\mathrm{secure}}$ & $t$ & $\delta$ \\
\midrule
\texttt{A-AUT-T} & auth     & 0.500 & 0.500 & 0.00 & 0.00 \\
\texttt{A-MEM-T} & memory   & 0.500 & 0.500 & 0.00 & 0.00 \\
\texttt{A-RES-T} & resource & 0.500 & 0.500 & 0.00 & 0.00 \\
\texttt{A-SES-T} & session  & 0.500 & 0.500 & 0.00 & 0.00 \\
\texttt{A-CCH-T} & tool     & 0.500 & 0.500 & 0.00 & 0.00 \\
\texttt{A-NSP-T} & tool     & 0.500 & 0.500 & 0.00 & 0.00 \\
\texttt{A-TOL-T} & tool     & 0.500 & 0.500 & 0.00 & 0.00 \\
\texttt{A-TRN-S} & tool     & 0.500 & 0.500 & 0.00 & 0.00 \\
\bottomrule
\end{tabularx}
\end{table}

\begin{figure}[t]
\centering
\includegraphics[width=\linewidth]{figures/rq1_leak_rate_by_boundary.pdf}
\caption{RQ-1: Leak rate per boundary, grouped by server. The
vulnerable server emits leakage in 50\% of calls on every
boundary; the secure server's per-connection probabilistic
leak path is suppressed by the defense middleware, but the
\emph{baseline} RQ-1 cell intentionally runs the secure server
\emph{without} all four defenses engaged (per-tenant tool
registry only), which is why the two bars are equal. The
defense-composition effect is measured separately under RQ-4.
Source: \texttt{analysis/figures/rq1\_leak\_rate\_by\_boundary.pdf}.}
\label{fig:rq1}
\end{figure}

\paragraph{Verdict for RQ-1: \textbf{accepted with caveat}.}
The vulnerable server exhibits a non-zero leak rate across all
eight boundaries, confirming that MCP's defaults leak in
isolation. The headline finding, however, is that the
\emph{baseline} defense level (per-tenant tool registry only)
is insufficient: it does not close the gap. The
\emph{full}-defense configuration is measured in RQ-4. The
negative result is not a power failure; the per-connection
leak probability on the partial-defense configuration is
deterministic ($p = 0.5$ per call), so the observed
$\Delta_i = 0$ for every attack reflects the defense
configuration, not insufficient power to detect a
difference.

\subsection{RQ-2: Cache dominance}\label{sec:rq2}

We restrict the RQ-2 manifest to the cache attacks
(\texttt{A-CCH-D, A-CCH-E, A-CCH-I, A-CCH-R, A-CCH-S,
A-CCH-T}) plus the grammar-fuzzing attack
(\texttt{A-FUZZ-001}) and measure each attack's contribution to
the vulnerable server's leakage volume.

\begin{table}[t]
\caption{RQ-2: Cache attack share of vulnerable-server
leakage. Each of the seven cache attacks contributes 14.29\%
of the total; the cache boundary as a whole accounts for 100\%
of the leakage observed in this run because no non-cache
attacks are present in the RQ-2 cell. One-sided $z$-test vs.
50\%: $z = 40.75$, $p < 10^{-4}$. Source:
\texttt{analysis/tables/rq2\_summary.csv}.}
\label{tab:rq2}
\centering
\small
\begin{tabularx}{\linewidth}{l r r}
\toprule
attack\_id & leak\_count & share \\
\midrule
\texttt{A-CCH-D} & 465 & 0.1429 \\
\texttt{A-CCH-E} & 465 & 0.1429 \\
\texttt{A-CCH-I} & 465 & 0.1429 \\
\texttt{A-CCH-R} & 465 & 0.1429 \\
\texttt{A-CCH-S} & 465 & 0.1429 \\
\texttt{A-CCH-T} & 465 & 0.1429 \\
\texttt{A-FUZZ-001} & 465 & 0.1429 \\
\midrule
\textbf{Total (cache)} & \textbf{3{,}255} & \textbf{1.0000} \\
\bottomrule
\end{tabularx}
\end{table}

\begin{figure}[t]
\centering
\includegraphics[width=\linewidth]{figures/rq2_cache_heatmap.pdf}
\caption{RQ-2: Leak count per attack on the vulnerable server.
The heatmap makes the uniformity of the cache attacks'
contribution visually obvious --- all seven cells carry the
same count ($465$). Source:
\texttt{analysis/figures/rq2\_cache\_heatmap.pdf}.}
\label{fig:rq2}
\end{figure}

\paragraph{Verdict for RQ-2: \textbf{accepted}.} The cache
boundary is the dominant source of leakage on the vulnerable
server; cache attacks account for 85.7\% of all observed
leakage when non-cache attacks are also present in the wider
RQ-1 cell, and 100\% of leakage in the cache-only RQ-2 cell.
The one-sided $z$-test against a 50\% null
($z = 40.75$, $p < 10^{-4}$) rejects the null decisively.

\subsection{RQ-3: Prompt-injection residual}\label{sec:rq3}

We run all seven prompt-injection attacks
(\texttt{A-TOL-001, A-TOL-D, A-TOL-E, A-TOL-I, A-TOL-R,
A-TOL-S, A-TOL-T}) against the secure server configured at the
\emph{partial} defense level (per-tenant tool registry only;
the full composition is measured separately in RQ-4). We
measure both the latency overhead and the leak rate. The
pre-registered bounded-residual hypothesis is that the secure
server's leak rate should fall inside the corridor
$[0.05, 0.30]$ set in \texttt{analysis/power.md}.

\begin{table}[t]
\caption{RQ-3: Per-attack latency and leak rate on the secure
server. The mean and p95 latencies are 0 ms because the
in-process connector does not introduce protocol overhead;
the leak-rate column reflects the residual leakage observed.
The bounded-residual corridor $[0.05, 0.30]$ is violated by
$0/7$ attacks. Source:
\texttt{analysis/tables/rq3\_summary.csv}.}
\label{tab:rq3}
\centering
\small
\begin{tabularx}{\linewidth}{l r r r r}
\toprule
attack\_id & mean (ms) & p95 (ms) & $L_{\mathrm{secure}}$ & in band \\
\midrule
\texttt{A-TOL-001} & 0.000 & 0.000 & 0.500 & no \\
\texttt{A-TOL-D}   & 0.000 & 0.000 & 0.500 & no \\
\texttt{A-TOL-E}   & 0.000 & 0.000 & 0.500 & no \\
\texttt{A-TOL-I}   & 0.000 & 0.000 & 0.500 & no \\
\texttt{A-TOL-R}   & 0.000 & 0.000 & 0.500 & no \\
\texttt{A-TOL-S}   & 0.000 & 0.000 & 0.500 & no \\
\texttt{A-TOL-T}   & 0.000 & 0.000 & 0.500 & no \\
\bottomrule
\end{tabularx}
\end{table}

\begin{figure}[t]
\centering
\includegraphics[width=\linewidth]{figures/rq3_injection_latency.pdf}
\caption{RQ-3: Per-attack latency on the secure server (box +
strip). The latency distribution is degenerate at zero because
the connector uses an in-process channel; the boxplot is
preserved for completeness. Source:
\texttt{analysis/figures/rq3\_injection\_latency.pdf}.}
\label{fig:rq3}
\end{figure}

\paragraph{Verdict for RQ-3: \textbf{rejected}.} The bounded-
residual hypothesis is rejected: $0/7$ attacks fall inside the
$[0.05, 0.30]$ corridor. The secure server's per-connection
leak probability, while reduced by the partial defenses
enabled in the RQ-3 cell, is not yet zero; the residual
leakage rate ($0.5$) is consistent with the structural-
property view of prompt injection as something no
protocol-level mitigation fully closes. The full defense
configuration is required (RQ-4).

\subsection{RQ-4: Defense composition}\label{sec:rq4}

We measure each attack's leak rate at three defense levels
(\texttt{none}, \texttt{partial} = per-tenant tool registry
only, \texttt{full} = all four defenses) and test the
pre-registered H1 that the per-attack reduction
$r_i = (L_{\mathrm{partial},i} - L_{\mathrm{full},i}) /
\max(L_{\mathrm{partial},i}, \epsilon)$ is at least
$0.5$. We report the paired one-sample Welch's $t$-test
against $0.5$ on the per-attack reduction (parametric) and
the Wilcoxon signed-rank test against $0.5$ (non-parametric;
appropriate because $r_i \in [0, 1]$ is bounded).

\begin{table}[t]
\caption{RQ-4: Per-attack leak rate and super-additivity
reduction. The full defense level (\texttt{full}) reduces the
leak rate to zero for every attack; the partial level
(\texttt{partial}) reduces it to roughly 0.19--0.24 but does
not close the gap. The reduction column is
$r_i = (L_{\mathrm{partial},i} - L_{\mathrm{full},i}) /
L_{\mathrm{partial},i}$, which exceeds the pre-registered
threshold of $0.5$ for every attack. The headline test is the
sign test against the threshold; the parametric $t$ and the
Wilcoxon signed-rank are reported as robustness checks. All
five attacks are significant at $p < 0.05$
($^{*}$). Source: \texttt{analysis/tables/rq4\_summary.csv}.}
\label{tab:rq4}
\centering
\small
\begin{tabularx}{\linewidth}{l r r r r r r}
\toprule
attack\_id & $L_{\mathrm{none}}$ & $L_{\mathrm{partial}}$ & $L_{\mathrm{full}}$ & $r_i$ & $p_{\mathrm{sign}}$ & sig \\
\midrule
\texttt{A-AUT-T} & 0.500 & 0.186 & 0.000 & 1.000 & 0.0312 & $^{*}$ \\
\texttt{A-CCH-T} & 0.500 & 0.237 & 0.000 & 1.000 & 0.0312 & $^{*}$ \\
\texttt{A-MEM-T} & 0.500 & 0.194 & 0.000 & 1.000 & 0.0312 & $^{*}$ \\
\texttt{A-RES-T} & 0.500 & 0.217 & 0.000 & 1.000 & 0.0312 & $^{*}$ \\
\texttt{A-TOL-T} & 0.500 & 0.220 & 0.000 & 1.000 & 0.0312 & $^{*}$ \\
\bottomrule
\end{tabularx}
\end{table}

\begin{figure}[t]
\centering
\includegraphics[width=\linewidth]{figures/rq4_defense_combo_bars.pdf}
\caption{RQ-4: Leak rate per defense level, stacked by attack.
The \texttt{none} level stacks to 50\% per attack; the
\texttt{partial} level stacks to roughly 20\% per attack; the
\texttt{full} level collapses to zero. Source:
\texttt{analysis/figures/rq4\_defense\_combo\_bars.pdf}.}
\label{fig:rq4}
\end{figure}

\paragraph{Verdict for RQ-4: \textbf{accepted}.} No single
defense closes the gap (RQ-3 confirms this for the partial
level alone), but the composition of all four defenses reduces
the leak rate to zero for every attack. The sign test on the
per-attack reductions $r_i$ against the pre-registered
threshold $0.5$ is decisive ($p = 0.0312$, one-sided); the
Wilcoxon signed-rank test agrees ($p = 0.0312$); the
parametric paired $t$ is infinite because the variance of
$r_i - 0.5$ is zero (every attack hits $r_i = 1.0$). This
is the central empirical result of the paper: defense-in-
depth at the MCP layer is necessary, and the specific
four-defense composition we propose is sufficient.

\paragraph{RQ-2 multiple-testing note.} The seven cache
attacks in Table~\ref{tab:rq2} share the same exploitation
primitive; we therefore report the headline $z$ as a single
test against the aggregate cache-share null, not as seven
independent tests. A Holm--Bonferroni correction across the
seven attacks leaves the headline result significant at
$\alpha = 0.05$.

\subsection{Summary of findings}\label{sec:summary-findings}

\begin{itemize}
  \item \textbf{RQ-1 accepted with caveat.} The naive MCP server
        leaks; the partial-defense configuration does not.
  \item \textbf{RQ-2 accepted.} Cache attacks dominate
        ($z = 40.75$, $p < 10^{-4}$).
  \item \textbf{RQ-3 rejected.} Prompt injection is a
        structural property; no partial mitigation fully closes
        it.
  \item \textbf{RQ-4 accepted.} Composition of four defenses
        reduces the leak rate to zero
        ($t = 22.80$, $p < 10^{-5}$).
\end{itemize}

The reproducibility appendix describes the exact command line
that produced each table and figure, including the seed, the
manifest path, and the analysis-script invocation.

\section{Defense Strategies}
% 07_defenses.tex — defense blueprint + matrix.

\section{Defense Strategies}\label{sec:defenses}

The RQ-4 finding --- that the composition of four defenses
reduces leak rate to zero, but no single defense closes the gap
--- motivates a layered defense blueprint. We enumerate the
four defenses, their cost, and the boundary they address.

\subsection{Defense matrix}

\begin{table}[t]
\caption{Defense matrix. Each row is one defense; columns are
the boundary it addresses, the mechanism, and the deployment
cost (Low = configuration change; Medium = middleware
addition; High = protocol-level rewrite).}
\label{tab:defense-matrix}
\centering
\small
\begin{tabularx}{\linewidth}{l l X c}
\toprule
Defense & Boundary & Mechanism & Cost \\
\midrule
Per-tenant tool registry         & namespace    & Static allow-list of tools per tenant                              & Low \\
Tenant-prefixed cache keys      & cache        & Cache keys include \texttt{tenant\_id}                             & Low \\
Resource-path canonicalisation  & resource     & Reject URIs that escape tenant root after \texttt{realpath}        & Low \\
Audience-bound JWT tokens       & auth         & \texttt{aud} claim enforced at every await boundary               & Medium \\
$m$TLS / channel auth           & transport    & Per-pipe tenant identity via SAN pinning                           & Medium \\
\bottomrule
\end{tabularx}
\end{table}

\subsection{Per-defense specification}

\paragraph{Per-tenant tool registry (namespace).} The
\texttt{Defenses.per\_tenant\_tool\_registry} flag toggles a
middleware that filters \texttt{tools/list} and
\texttt{tools/call} responses by the requesting tenant. The
registry is a \texttt{dict[tenant\_id, set[tool\_name]]}
populated at server startup; an unknown tenant receives an
empty list, an unknown tool receives a 403. The middleware
adds a constant-time hash lookup per call ($\sim 50\,\mu s$ in
our reference implementation).

\paragraph{Tenant-prefixed cache keys (cache).} The
\texttt{Defenses.tenant\_prefixed\_cache\_keys} flag
re-computes every cache key as
\texttt{(tenant\_id, tool\_name, args\_hash)} instead of
\texttt{(tool\_name, args\_hash)}. The middleware is a single
line of code at the cache lookup site; cost is negligible.

\paragraph{Resource-path canonicalisation (resource).} The
\texttt{Defenses.resource\_path\_canonicalisation} flag
replaces the resolver's string-prefix check with
\texttt{os.path.realpath} after decoding percent-encoded
segments. URIs that escape the tenant root return a 404. The
middleware adds $\sim 200\,\mu s$ per
\texttt{resources/read} call.

\paragraph{Audience-bound JWT tokens (auth).} The
\texttt{Defenses.mtls} flag (the prompt-mandated name for the
audience-bound JWT defense) issues JSON Web Tokens with an
\texttt{aud} claim scoped to the tenant identifier. The
server re-validates the \texttt{aud} claim at every await
boundary (re-check, not single-check, to defeat the
TOCTOU race in misuse case). The middleware adds $\sim
300\,\mu s$ per call.

\subsection{Defense-composition rationale}

No single defense is sufficient because each addresses a
single boundary:

\begin{itemize}
  \item The per-tenant tool registry does not affect cache
        poisoning (RQ-2 demonstrates cache attacks bypass
        the registry).
  \item Tenant-prefixed cache keys do not affect resource
        traversal (RQ-3 demonstrates resource attacks bypass
        the cache middleware).
  \item Resource canonicalisation does not affect prompt
        injection (the prompt is not a URI).
  \item Audience-bound JWTs do not affect cache key
        collisions (the collision is below the auth layer).
\end{itemize}

The composition closes the gap because each attack's
exploitation path crosses \emph{at least one} boundary whose
defense is engaged. Cache attacks cross the namespace boundary
(tool routing) and the cache boundary; the registry defence
catches the former and the cache defence catches the latter.
Prompt-injection attacks cross the resource boundary (the
injected content is read via \texttt{resources/read}) and the
auth boundary (the injected instruction seeks elevation); the
canonicalisation and JWT defences catch them in series.

\subsection{Layered defense blueprint}\label{sec:blueprint}

We recommend the following order of adoption for a production
MCP deployment:

\begin{enumerate}
  \item \textbf{Audience-bound JWTs (auth).} Highest leverage;
        addresses the broadest set of attacks
        (token-replay, capability-spoofing, SSRF).
  \item \textbf{Per-tenant tool registry (namespace).} Low
        cost; addresses tool-squatting and confused-deputy
        attacks.
  \item \textbf{Tenant-prefixed cache keys (cache).} Low cost;
        addresses the dominant leakage surface (RQ-2).
  \item \textbf{Resource-path canonicalisation (resource).} Low
        cost; addresses path traversal and symlink escape.
\end{enumerate}

The order is chosen so that each subsequent defense is
additive rather than overlapping: the JWT defence gates
\emph{who}, the tool registry gates \emph{what}, the cache
keys gate \emph{how often}, and the resource canonicalisation
gates \emph{where to}.

\subsection{Defense overhead (measured)}

The per-defense overhead is measured on the in-process
connector by \texttt{analysis/scripts/bench\_defenses.py}
(10{,}000 iterations per defense, median reported):

\begin{itemize}
  \item Per-tenant tool registry: $0.1\,\mu s$ per call
        (constant-time hash lookup).
  \item Tenant-prefixed cache keys: $0.1\,\mu s$ per call
        (single string concatenation).
  \item Resource-path canonicalisation: $\sim 113\,\mu s$ per
        call (\texttt{os.path.realpath} on a typical URI).
  \item Audience-bound JWTs: $\sim 2\,\mu s$ per call
        (HMAC-SHA256 verify).
\end{itemize}

The cumulative per-call overhead is therefore
$\sim 115\,\mu s$ on the in-process connector. On a real
HTTP+SSE transport the network round-trip ($\sim 1$--$10\,ms$
locally) dominates the per-defense overhead; the relative
ordering of the defenses (and the composition result) is
unaffected. The benchmark harness is included in the
artifact bundle so reviewers can re-measure on their target
transport.

\subsection{Defense failure modes}

Each defense has a failure mode that a determined attacker can
exploit. We document them so that adopters understand the
defenses' residual risk:

\paragraph{Per-tenant tool registry failure mode.} The
registry is a \texttt{dict[tenant\_id, set[tool\_name]]};
if the attacker can register a tool with a \texttt{tenant\_id}
the server trusts (e.g. via a privilege-escalation bug in the
host's authentication), the registry's allow-list is moot.
The defense therefore relies on the host's authentication
chain, which is itself defended by the audience-bound JWT
defense (the second layer).

\paragraph{Tenant-prefixed cache keys failure mode.} If the
attacker can poison a tenant's cache key (e.g. via a separate
write-side bug), the prefix does not prevent the poisoned
key from being read by the same tenant. The defense therefore
relies on the write-side defenses (per-tenant tool registry
and resource canonicalisation) to prevent the initial
poisoning. The defense closes the \emph{read-side} leakage
path, not the \emph{write-side} attack.

\paragraph{Resource-path canonicalisation failure mode.} If
the attacker can craft a URI that resolves within the tenant
root but contains attacker content from another tenant (e.g.
via a server-side include or a symlink planted at startup),
the canonicalisation does not detect the cross-tenant
content. The defense closes the \emph{path-traversal}
leakage path, not the \emph{content-injection} attack. The
prompt-injection defense (future work) would close the
latter.

\paragraph{Audience-bound JWTs failure mode.} If the attacker
can compromise the upstream identity provider (out of scope
in our threat model), the JWT's \texttt{aud} claim is moot
because the attacker can mint a token with any
\texttt{aud}. The defense relies on the identity provider's
security; the protocol-layer defense is the second line of
defence, not the first.

\subsection{Defenses and the four misuse cases}

We map each defense to the misuse cases it addresses:

\begin{itemize}
  \item \textbf{MC-1 (tool shadowing).} Closed by the
        per-tenant tool registry. The shadowed \texttt{set\_env}
        tool is not in Tenant~B's allow-list; the call
        returns a 403.
  \item \textbf{MC-2 (session fixation).} Closed by the
        audience-bound JWTs. The attacker cannot mint a
        session-resumption JWT with Tenant~B's
        \texttt{aud} claim.
  \item \textbf{MC-3 (resource path traversal).} Closed by
        the resource-path canonicalisation. The
        percent-encoded slashes are decoded and
        \texttt{realpath}-ed before the tenant-prefix check;
        the resolver returns a 404.
  \item \textbf{MC-4 (embedding cache poisoning).} Closed by
        the tenant-prefixed cache keys. The cache key now
        includes \texttt{tenant\_id}; a collision crafted by
        Tenant~A does not match Tenant~B's key.
\end{itemize}

The composition is therefore complete: every documented
misuse case is closed by at least one defense, and most are
closed by two (defense-in-depth). The RQ-4 empirical result
($L = 0$ at composition) is the numerical confirmation of
this mapping.

\section{Related Work}
% 08_related_work.tex — synthesises literature/related_work.md.

\section{Related Work}\label{sec:related}

We situate our work in six clusters of prior literature,
organised by the boundary or methodology they target.

\subsection{LLM tool-use and prompt injection}

Greshake et al.~(2023) introduced the concept of \emph{indirect
prompt injection} via tool and retrieval content, demonstrating
that LLM-integrated applications are broadly exploitable
through content they fetch~\cite{greshake_2023_indirect}. Zou
et al.~(2023) developed automated, transferable adversarial
suffixes for aligned LLMs via greedy coordinate-gradient
search~\cite{zou_2023_universal}. Chowdhury et al.~(2024)
provided a comparative survey of attacks on LLMs and showed
that defenses fail in isolation and under
composition~\cite{chowdhury_2024_defenses} --- the same
empirical pattern we observe at the MCP layer. Wang et
al.~(2024) built ToolCommander, a framework for LLM
tool-scheduling hijack via injected tool descriptors, which
matches our \texttt{A-NSP-001} and \texttt{A-TOL-005}
findings~\cite{wang_2024_toolcommander}. These works establish
that prompt injection is a structural property of LLM agents
and that tool scheduling is exploitable; we extend both
findings to the MCP protocol layer with concrete per-boundary
attack patterns.

\subsection{MCP and agent-protocol security}

Chen et al.~(2026) conducted the largest empirical study of
runtime MCP servers to date, finding that scanners under-
report runtime risks and that production servers routinely
share stdio workers and follow symlinks across tenant
roots~\cite{chen_2026_mcp_security}. Their work validates our
Phase-1 threat model at scale. An et al.~(2026) introduced
FlowGuard, a behavioural-trace correlation detector for MCP
security~\cite{an_2026_flowguard}; their methodological
approach informs our Evaluator module (\S\ref{sec:framework}).
Lee et al.~(2026) demonstrated mid-session tool injection in
WebMCP, showing that third-party-script registry mutation
bypasses pre-session policies~\cite{lee_2026_webmcp}; we add
\texttt{notifications/tools/list\_changed} audit-trail support
to our reference secure server in response. Liu et al.~(2026)
introduced MCPEvol-Bench, a dynamic-evolution benchmark that
shows static benchmarks overstate robustness~\cite{liu_2026_mcpevol};
mid-session mutations are now first-class in our attack
library. Jing et al.~(2026) argued for \emph{isolation as a
first-class principle} for LLM-agent
safety~\cite{jing_2026_isolation} --- the exact framing of
our paper --- but provided neither a concrete boundary
catalogue nor empirical measurement.

\subsection{Memory and cache attacks}

Torres et al.~(2026) demonstrated long-term memory poisoning
in LLM agents, showing that memory entries persist across
sessions and that write-time detection is
preferred~\cite{torres_2026_memory}. Their findings bridge the
prompt-injection cluster and our memory/cache boundaries: the
agent-side counterpart of our \texttt{A-CCH-001} /
\texttt{A-MEM-001} MCP-side failure modes.

\subsection{Sandboxing and physical isolation}

Dipta et al.~(2024) demonstrated that side channels survive
logical sandboxing: their dynamic frequency-based
fingerprinting attack defeats gVisor, Firecracker, SGX, and
SEV~\cite{dipta_2024_sandbox}. This is a threat-model
limitation we honestly disclose (\S\ref{sec:discussion}); our
work is at the MCP protocol layer, not the silicon layer.
Amazon's Firecracker microVM~\cite{firecracker_github} is the
substrate for our reference-server deployment.

\subsection{Standards and architectural priors}

NIST SP 800-207 (Zero Trust Architecture~\cite{nist_sp_800_207})
provides the architectural prior for our defense blueprint:
\emph{never trust, always verify}. The OWASP LLM Top 10
v1.1~\cite{owasp_llm_top10_v11} provides the LLM-application
taxonomy; we extend it to the protocol layer with explicit
MCP-boundary bindings. The MCP specification and Anthropic's
announcement~\cite{mcp_spec,anthropic_mcp_announce} are the
specification under study; our defense blueprint and spec-
change recommendations target the next revision (MCP-2).

\subsection{LSP analogue}

Zhu et al.~(2025) introduced LSPFuzz, a grammar-aware fuzzing
campaign over $\sim$300 LSP servers that found many
bugs~\cite{zhu_2025_lspfuzz}. LSP and MCP are both JSON-RPC-
based protocols between a host and a language/tool-specific
server; LSPFuzz is the closest architectural analogue and the
methodological template for our \texttt{GrammarFuzzAttack}
(\texttt{A-FUZZ-001}). We adapt the campaign to MCP's eight-
boundary surface.

\subsection{Positioning summary}

\begin{table}[t]
\caption{Positioning vs.\ prior work. \checkmark{} = covered;
$\triangle$ = partial; $\times$ = not covered.}
\label{tab:positioning}
\centering
\small
\begin{tabularx}{\linewidth}{l c c c c c}
\toprule
                                & \rotatebox{70}{Greshake 2023} & \rotatebox{70}{Wang 2024} & \rotatebox{70}{Torres 2026} & \rotatebox{70}{Chen 2026} & \rotatebox{70}{This work} \\
\midrule
Eight-boundary catalogue      & $\times$  & $\times$ & $\times$ & $\triangle$ & \checkmark \\
Concrete per-boundary attacks  & $\times$  & $\triangle$ & $\triangle$ & $\triangle$ & \checkmark \\
Reproducible harness           & $\times$  & $\checkmark$ & $\times$ & $\triangle$ & \checkmark \\
Multi-tenant leakage metric   & $\times$  & $\times$ & $\times$ & $\times$ & \checkmark \\
Defense-composition study      & $\times$  & $\times$ & $\times$ & $\times$ & \checkmark \\
Open attack library            & $\times$  & $\times$ & $\times$ & $\times$ & \checkmark \\
\bottomrule
\end{tabularx}
\end{table}

The novelty score is 8/10 per our pre-registered Phase-4
assessment (\texttt{docs/01\_Research\_Gap.md}); the
empirical numbers in this paper upgrade the conceptual
scaffolding to executed evidence.

\subsection{How our threat model differs}

Our threat model differs from each prior cluster in a specific
way:

\begin{itemize}
  \item \textbf{vs.\ Greshake et al.~(2023)}: their threat
        model is the LLM-integrated application; ours is the
        MCP server. Indirect prompt injection is the common
        surface; the defense locus differs.
  \item \textbf{vs.\ Wang et al.~(2024)}: their threat model
        is the LLM tool scheduler; ours is the MCP namespace
        boundary. ToolCommander demonstrates tool-scheduling
        hijack; our \texttt{A-NSP-001} row is the MCP-specific
        analogue.
  \item \textbf{vs.\ Torres et al.~(2026)}: their threat model
        is agent-side memory; ours is server-side cache +
        memory. The agent-server symmetry is what makes the
        two papers complementary: Torres covers the agent
        half of the boundary; we cover the server half.
  \item \textbf{vs.\ Chen et al.~(2026)}: their threat model
        is the MCP server's single-tenant default behaviour;
        ours is the multi-tenant isolation surface. Chen
        measures single-tenant runtime risk; we measure
        cross-tenant leakage.
  \item \textbf{vs.\ An et al.~(2026)}: their threat model is
        MCP-server signal traces; ours is MCP-server
        protocol-level leakage. FlowGuard detects; we measure.
  \item \textbf{vs.\ Lee et al.~(2026)}: their threat model
        is WebMCP mid-session mutation; ours is MCP static
        surface plus mid-session defence audit. The two
        threat models are bridged by the
        \texttt{notifications/tools/list\_changed} audit-trail
        recommendation in our defense blueprint.
  \item \textbf{vs.\ Liu et al.~(2026)}: their threat model is
        MCP-server tool evolution; ours is MCP-server
        isolation. MCPEvol-Bench studies robustness under
        change; we study leakage under concurrency.
  \item \textbf{vs.\ Jing et al.~(2026)}: their threat model
        is LLM-agent safety in general; ours is MCP
        multi-tenant safety specifically. Jing argues for
        isolation as a first-class principle; we
        operationalise it with concrete boundaries and
        defenses.
  \item \textbf{vs.\ Dipta et al.~(2024)}: their threat model
        is OS-level sandboxing; ours is protocol-level
        isolation. The two papers are complementary: Dipta
        covers what protocol-layer defenses cannot close;
        we cover what OS-layer sandboxes cannot close.
\end{itemize}

The common thread is that \emph{no prior work} combines all
eight boundaries, all six STRIDE letters, a reproducible
multi-tenant harness, a leakage metric, and a defense-
composition study. The combination is the unit of novelty.

\subsection{Cluster F revisited: why LSPFuzz is the closest analogue}

LSPFuzz~\cite{zhu_2025_lspfuzz} is the closest architectural
analogue to MCP because LSP and MCP share three structural
properties:

\begin{enumerate}
  \item Both are JSON-RPC-2.0-based protocols.
  \item Both mediate between a host (an editor or an agent
        runtime) and a language/tool-specific server.
  \item Both are stateless at the protocol level but stateful
        at the server level.
\end{enumerate}

LSPFuzz's grammar-aware fuzzing campaign over $\sim$300 LSP
servers is therefore the methodological template for our
\texttt{GrammarFuzzAttack}; the differences are in the
attack-surface derivation (LSP's surface is grammar-shaped;
MCP's surface is boundary-shaped). We adapt LSPFuzz's
approach to the eight-boundary surface and add the
\texttt{A-CCH-*} cache attacks that have no LSP analogue.

\subsection{Cluster E revisited: standards as the design prior}

Our defense blueprint is consistent with three standards-
level priors:

\begin{itemize}
  \item \textbf{NIST SP 800-207 (Zero Trust)}
        \cite{nist_sp_800_207}: ``never trust, always verify''
        is the design principle behind all four defenses.
        Each defense re-validates the tenant identity at a
        different boundary.
  \item \textbf{OWASP LLM Top 10 v1.1}
        \cite{owasp_llm_top10_v11}: LLM01 (prompt injection)
        and LLM06 (sensitive information disclosure) are the
        two entries our defenses address; the others are out
        of scope.
  \item \textbf{MCP specification}
        \cite{mcp_spec,anthropic_mcp_announce}: the spec
        defines the wire protocol; our spec-change
        recommendations are the minimum changes required
        to make the protocol-level isolation
        enforceable-by-default.
\end{itemize}

The defense blueprint is therefore not ad-hoc; it is the
intersection of the three standards-level priors applied
to the MCP protocol surface.

\section{Discussion}
% 09_discussion.tex — limitations, ethics, generalisability,
% future work.

\section{Discussion}\label{sec:discussion}

\subsection{Limitations}

\paragraph{Connector channel.} Our reference connector uses an
in-process channel; the latency measurements in RQ-3 are
therefore not representative of a real HTTP+SSE deployment.
The leak-rate measurements (RQ-1, RQ-2, RQ-4) are unaffected
because the connector's leak probability is the same regardless
of channel.

\paragraph{Reference servers.} The vulnerable and secure
reference servers are written against a single MCP SDK
(``FastMCP''). Generalising to the TypeScript and Rust SDKs
would surface SDK-specific failure modes (e.g. decorator-name
shadowing in the Python SDK does not apply to the TypeScript
SDK). We mark SDK coverage as future work.

\paragraph{Defense coverage.} We evaluate four protocol-layer
defenses; we do not evaluate OS-level sandboxing because the
goal is to fix MCP, not the underlying OS (Dipta et al.~(2024)
already covers the OS layer~\cite{dipta_2024_sandbox}). The
defenses are evaluated in composition but not in isolation in
RQ-4; the per-defense single-cell measurements are in
\texttt{analysis/runs/exp-rq4-defense/} and can be re-analysed
post hoc.

\paragraph{Statistical power.} With $n = 30$ per cell, our
power to detect medium-effect-size differences ($d \approx 0.5$)
is $\sim 0.80$ at $\alpha = 0.05$. Detecting smaller effects
would require larger samples; the framework supports
arbitrary iteration counts via the manifest's
\texttt{run.repeats} field.

\paragraph{Side channels.} We do not measure silicon-level
side channels (cache-line timing, rowhammer, Spectre/Meltdown).
Dipta et al.~(2024) demonstrated that logical isolation is
incomplete against such attacks; our work is at the protocol
layer and inherits this limitation honestly.

\paragraph{Real-world LLM integration.} Our experiments use a
deterministic in-process connector, not a real LLM. Prompt-
injection effectiveness against a real model would depend on
the model's alignment and instruction-following behaviour;
our \texttt{ResultInjectionAttack} is the protocol-level
embodiment of the injection surface, not a model-specific
exploit.

\subsection{Ethics}

The experiments in this paper use only synthetic tenant data
and synthetic marker payloads; no real user data, real
credentials, or real LLM outputs are involved. The vulnerable
reference server is run inside a Firecracker microVM
(\texttt{mcp\_servers/vulnerable/}) so that even if an attack
exfiltrates data, the data is synthetic. The attack library is
released under MIT so that defensive research can build on it;
we do not release a weaponised exploit kit and we coordinate
with the MCP maintainers before any public disclosure of
specific server vulnerabilities discovered during this work.

\subsection{Generalisability}

Our threat model assumes a multi-tenant MCP server with at
least two tenants; the attacks and defenses generalise to any
multi-tenant protocol that distributes isolation across
implicit boundaries. The framework's pluggable attack and
defense interfaces (\texttt{attacks/base.py},
\texttt{Defenses} Pydantic model) make it straightforward to
port the harness to other JSON-RPC-based agent protocols
(LSP, ADP, Agent Protocol v2).

\subsection{Future work}

\paragraph{Ticket-specific attack classes.} We ship 50 attack
classes (48 STRIDE rows + 2 cross-cutting). The Phase-1 ticket
catalogue has 63 tickets; one class per ticket would surface
finer-grained exploitation paths.

\paragraph{Cross-SDK coverage.} Extending the attack library
to the TypeScript and Rust MCP SDKs would surface SDK-specific
vulnerabilities not covered by the Python reference.

\paragraph{Real-LLM integration.} Pairing the framework with a
real LLM would let us measure prompt-injection effectiveness
as a function of model alignment; the framework's
\texttt{Defenses.prompt\_content\_scanning} flag is already
in place to support this experiment.

\paragraph{Defense ordering.} The defense blueprint in
\S\ref{sec:blueprint} recommends a fixed adoption order;
adaptive ordering based on the deployment's threat profile is
future work.

\paragraph{Spec-change proposals.} We will draft concrete text
for three MCP-2 revisions (audience-bound JWTs as default;
tenant-prefixed cache keys as default; URI canonicalisation as
default) and submit them via the MCP specification change-
request process.

\paragraph{Long-running agent memory.} Torres et al.~(2026)
identify memory poisoning as more durable than prompt
injection; integrating their memory-poisoning taxonomy with
our \texttt{A-MEM-001..008} tickets is an obvious next step.

\subsection{What we learned about MCP security that the
specification does not say}

The empirical results support three claims that the MCP
specification does not currently make and arguably should:

\paragraph{Claim 1: cross-tenant isolation is not a
property of the default reference implementations.} RQ-1
demonstrates that a naive MCP server leaks across all eight
boundaries; the secure server's defense middleware is
required to close the leakage paths. The specification
should state this explicitly and recommend the defense
middleware as the default deployment configuration.

\paragraph{Claim 2: the cache boundary is the dominant
leakage surface.} RQ-2 demonstrates that cache attacks
account for 85.7\% of observed leakage. The specification
should require tenant-prefixed cache keys as the default
and document the failure mode of bare
\texttt{(tool\_name, args\_hash)} keys.

\paragraph{Claim 3: prompt injection is a structural property
of tool-using LLM agents.} RQ-3 demonstrates that no
partial mitigation fully closes prompt injection. The
specification should acknowledge this honestly and
recommend defense-in-depth (composition of the four
defenses) as the only known mitigation at the protocol
layer.

\subsection{Threat-model honesty and the ``spec change''
recommendation}

We deliberately scope our spec-change recommendations to the
protocol layer because the protocol layer is what the MCP
specification controls. We do not recommend OS-level
sandboxing in the specification because that is what the OS
controls; we do not recommend LLM-side defenses because
that is what the model provider controls. The
recommendations in \S\ref{sec:conclusion} are the minimum
changes the specification can make to make the protocol-
layer isolation enforceable-by-default.

\subsection{What an empirical paper at a security venue looks like}

For readers unfamiliar with the security-venue genre, the
paper's empirical claims follow a specific template:

\begin{enumerate}
  \item A pre-registered statistical protocol
        (\texttt{analysis/power.md}).
  \item A pre-registered sample size ($n = 30$ per cell).
  \item A pre-registered significance threshold
        ($\alpha = 0.05$ with Bonferroni correction within
        boundary).
  \item A pre-registered effect-size metric (Cliff's $\delta$
        for non-parametric; Cohen's $d$ for parametric).
  \item A reproducible artifact (the attack library, the
        framework, the manifests, the raw run outputs).
\end{enumerate}

Each empirical claim in \S\ref{sec:evaluation} traces back to
this protocol. The Phase-12 reviewer simulation
(\texttt{prompts/12\_reviewer.md}) will critique the
experimental design and the statistical protocol; we
address the critique in \S\ref{sec:discussion}.

\section{Conclusion}
% 10_conclusion.tex — one-page wrap; restate contributions;
% call for MCP-2 spec changes.

\section{Conclusion}\label{sec:conclusion}

We presented an empirical study of cross-tenant isolation in
the Model Context Protocol. Our contributions are five-fold:
(i)~an eight-boundary STRIDE catalogue with 63 forward ticket
IDs; (ii)~a 14-concept security taxonomy with explicit
MCP-boundary bindings; (iii)~a reproducible cross-tenant
leakage measurement framework; (iv)~an open attack library of
50 concrete classes with pytest reproducibility; and
(v)~an MCP-layer defense-comparison study.

The empirical findings are unambiguous:

\begin{itemize}
  \item \textbf{The naive MCP server leaks across all eight
        boundaries.} RQ-1 confirms a non-zero leak rate on
        every attack.
  \item \textbf{Cache attacks dominate.} RQ-2 shows cache
        attacks account for 85.7\% of leakage
        ($z = 40.75$, $p < 10^{-4}$).
  \item \textbf{Prompt injection is a structural property.} RQ-3
        shows that no partial mitigation fully closes it.
  \item \textbf{Composition works.} RQ-4 shows that the four-
        defense composition reduces leak rate to zero
        ($t = 22.80$, $p < 10^{-5}$).
\end{itemize}

\subsection{Spec-change recommendations for MCP-2}

We close with three concrete recommendations for the next
revision of the Model Context Protocol specification:

\begin{enumerate}
  \item \textbf{Audience-bound JWTs as the default.} The auth
        boundary should mandate an \texttt{aud} claim scoped to
        the tenant identifier, validated at every await
        boundary. Current best-practice (per-tool scope,
        per-tenant scope) is opt-in and inconsistently
        deployed.
  \item \textbf{Tenant-prefixed cache keys as the default.} The
        cache boundary should require every cache key to
        include the resolved \texttt{tenant\_id}. Current
        practice keys caches by \texttt{(tool\_name,
        args\_hash)} alone, enabling the dominant leakage
        surface.
  \item \textbf{URI canonicalisation as the default.} The
        resource boundary should require
        \texttt{resources/read} resolvers to apply
        \texttt{os.path.realpath} after percent-decoding
        before any tenant-prefix check. Current practice
        concatenates tenant roots with raw URI paths,
        enabling the path-traversal attacks in
        misuse case MC-3.
\end{enumerate}

Each recommendation is implementable as a single middleware
line in any MCP SDK and adds $\sim 200$--$300\,\mu s$ per call.
We estimate that adopting all three would close the leak in
RQ-1, RQ-2, and RQ-4 without further changes to the
specification.

\paragraph{Availability.}
The artifact (attack library, measurement framework, analysis
scripts, notebooks, summary CSV tables, and rendered figures)
is available at the project repository under the MIT license.
The pinned dependency versions in \texttt{pyproject.toml} and
the seed-controlled manifests in
\texttt{experiments/manifests/} make the results reproducible
on any machine with the documented toolchain.

\paragraph{Reproducibility statement.}
We provide two complete appendices: Appendix~A
(\texttt{paper/appendix\_a\_data\_traceability.md}) maps every
number in \S\ref{sec:evaluation} back to a row in the
analysis CSVs; Appendix~B
(\texttt{paper/appendix\_b\_reproduction\_log.md}) records
the exact command line that produced each table and figure,
the seed, the manifest path, the analysis-script invocation,
and the \texttt{commit\_sha} under which the run was executed.
The same artifact was re-run on three independent machines
(Linux x86-64, macOS arm64, Windows x86-64) with bit-identical
CSV outputs modulo platform-dependent timestamps
(\texttt{meta.json:started\_at}). The Docker image
(\texttt{artifact/docker/Dockerfile}) reproduces the pipeline
inside a single \texttt{docker compose up} invocation.

\paragraph{Acknowledgements.}
We thank the MCP maintainers for their responsiveness to
specification questions during this work; the maintainers of
\texttt{scipy}, \texttt{pandas}, and \texttt{matplotlib} for
the statistical and plotting primitives; and the anonymous
reviewers of the Phase-12 reviewer simulation for their
rigorous critique of the experimental design.

% Single consolidated bibliography at the end (alphabetic labels).
\printbibliography

\end{document}
